\documentclass{subfiles}
\begin{document}
    We start by choosing a prime number \(p\) such that all possible key fall in the range 
        \((0, p - 1)\). Consider the group \(\zset_{p}\) and \(\zset_{p}^{*}\).
        Since \(p\) is prime we can solve equations modulo \(p\). 
    
    For any \(a \in \zset_{p}^{*}\) and \(b \in \zset_{p}\), 
        let \(h_{ab}\) be the  function defined as follow:
        \[
            h_{ab}(k) =  ((ak + b) \mod p) \mod m
        \]
        and let \(\H_{pm}\) be the class of all such functions.
    
    \begin{theorem}
        The family of hash functions \(\H_{pm}\) is universal.
    \end{theorem}
    \begin{proof*}
        For any two distinct keys \(k, l \in \zset_{p}\), and an hash function \(h_{ab}\) we have 
            \[\begin{gathered}
                h_{ab}(k) = r = ((ak + b) \mod p) \mod m \\ 
                h_{ab}(l) = s = ((al + b) \mod p) \mod m
            \end{gathered}\]
            Some simple algebra shows that \(r \ne s\). Hence,
            at the modulo \(p\) level no collision can occur.

        Therefore, the probability that two keys hash to the same value coicide with
            the probability that \(r \equiv s \mod m\). Now, if \(r\) and \(s\) are chosen at random 
            from \(\zset_{p}\), for a fixed value of \(r\),
            the number of values for \(s\) such that \(s \ne p\) and \(s \equiv p \mod m\) is at most
            \[\begin{aligned}
                \ceil{\sfrac{p}{m}} - 1 & \le (\sfrac{(p + m - 1)}{m}) - 1 \\ 
                & = \sfrac{(p - 1)}{m}.
            \end{aligned}\]
            Meaning that the probability that \(r \equiv s \mod m\) is at most \(\sfrac{\sfrac{(p - 1)}{m}}{(p - 1)} = \sfrac{1}{m}\).
            Hence, \(\Prob{h_{ab}(k) = h_{ab}(l)} \le \sfrac{1}{m}\), and thus \(\H_{pm}\) is universal.
    \end{proof*}
\end{document}
